\subsection{整体超 Tor 的结合性谱序列}
\source{\texttt{EGAIII.2.pdf}，PDF第31--32页；原书第32--33页}
\egap{6.8.1}
采用（6.7.11）的假设与记号，特别地，各 $\mathscr P^{(i)}_\bullet$ 下有界。设给定指标集 $I$ 的一个分割 $(I_j)_{j\in J}$。要在整体超 Tor
\[
 \lTor_*^S\bigl((f_i)_{i\in I};(\mathscr P^{(i)}_\bullet)_{i\in I}\bigr)
\]
与各“部分”超 Tor
\[
 \mathscr A^{(j)}_\bullet
 =\lTor_*^S\bigl((f_i)_{i\in I_j};(\mathscr P^{(i)}_\bullet)_{i\in I_j}\bigr)
\]
之间给出一种“结合性”关系。为简化书写，只考虑 $I=\{1,\ldots,m\}$，且分割由两个区间 $\{1,\ldots,r\}$ 与 $\{r+1,\ldots,m\}$ 组成的情形。

\egas{6.8.2}{命题}
存在典范的双正则谱函子，称为“结合性谱函子”，记作
\[
 {}^{(r)}\spE(f_1,\ldots,f_m;
                  \mathscr P^{(1)}_\bullet,\ldots,\mathscr P^{(m)}_\bullet)
\]
（也可在复形记号中略去下标点），其终值为
\[\lTor_*^S(f_1,\ldots,f_m;\mathscr P^{(1)}_\bullet,\ldots,\mathscr P^{(m)}_\bullet),\]且 $E^2$ 项为
\[
\begin{aligned}
 {}^{(r)}\spE^2_{pq}
 =\bigoplus_{q_1+q_2=q}\lTor_p^S\bigl(&
   \lTor_{q_1}^S(f_1,\ldots,f_r;
                      \mathscr P^{(1)}_\bullet,\ldots,\mathscr P^{(r)}_\bullet),\\
 &\lTor_{q_2}^S(f_{r+1},\ldots,f_m;
                      \mathscr P^{(r+1)}_\bullet,\ldots,\mathscr P^{(m)}_\bullet)\bigr).
\end{aligned}
\]
在此，已将 $Y$ 典范地等同于 $Z^{(1)}\times_SZ^{(2)}$，其中
\[
 Z^{(1)}=Y^{(1)}\times_S\cdots\times_SY^{(r)},\qquad
 Z^{(2)}=Y^{(r+1)}\times_S\cdots\times_SY^{(m)}.
\]
我们只处理 $S$ 和各 $Y^{(i)}$ 仿射的情形；用（6.7.1）、（6.7.2）的方法从这一特例过渡到一般情形，没有困难，论证细节留给读者。因此证明下面的推论即可。

\egas{6.8.3}{推论}
设 $A$ 为环，$S=\Spec A$，$X^{(i)}$（$1\leq i\leq m$）为拟紧 $S$-概形；对每个 $i$，设 $\mathscr P^{(i)}_\bullet$ 为下有界拟凝聚 $\OO_{X^{(i)}}$-模复形。存在一个典范双正则谱函子，其终值为
\[
 \bTor_*^S(X^{(1)},\ldots,X^{(m)};
                  \mathscr P^{(1)}_\bullet,\ldots,\mathscr P^{(m)}_\bullet),
\]
且 $E^2$ 项为
\[
\begin{aligned}
 {}^{(r)}E^2_{pq}=\bigoplus_{q_1+q_2=q}\Tor_p^A\bigl(&
   \bTor_{q_1}^S(X^{(1)},\ldots,X^{(r)};
                    \mathscr P^{(1)}_\bullet,\ldots,\mathscr P^{(r)}_\bullet),\\
 &\bTor_{q_2}^S(X^{(r+1)},\ldots,X^{(m)};
                    \mathscr P^{(r+1)}_\bullet,\ldots,\mathscr P^{(m)}_\bullet)\bigr).
\end{aligned}
\]
\begin{proof}
按（6.6.2）的定义，计算所讨论的超 Tor 时，对每个 $i$ 取 $X^{(i)}$ 的有限仿射开覆盖 $\mathfrak U^{(i)}$，考虑双复形
$L^{(i)}_{\bullet\bullet}=C^\bullet(\mathfrak U^{(i)},\mathscr P^{(i)}_\bullet)$，并取它的投射 Cartan--Eilenberg 解消 $M^{(i)}_{\bullet\bullet\bullet}$，这里解消按（0，11.7.1）理解。形成这些三重复形的张量积
\[
 M_{\bullet\bullet\bullet}=\bigotimes_{i=1}^mM^{(i)}_{\bullet\bullet\bullet},
\]
然后取其同调。把 $M$ 看作单复形 $N_\bullet$，它是下列两个单复形的张量积：
\[
 N'_\bullet=\bigotimes_{i=1}^rM^{(i)}_{\bullet\bullet\bullet},\qquad
 N''_\bullet=\bigotimes_{i=r+1}^mM^{(i)}_{\bullet\bullet\bullet};
\]
其中各单复形按相应总次数之和分次。$N'_\bullet,N''_\bullet$ 中各 $A$-模均投射，故（6.5.9）给出
$H_*(M)=\Tor_*^A(N'_\bullet,N''_\bullet)$。所求谱序列就是将（6.3.2.2）用于 $N'_\bullet,N''_\bullet$；上述构造用于两个部分乘积时，给出这两个复形的同调模的所需解释。最后，由于各 $\mathscr P^{(i)}_\bullet$ 下有界，各 $M^{(i)}$ 也下有界，因而 $N'_\bullet,N''_\bullet$ 下有界；所需的正则性由（6.3.2）得到。
\end{proof}
